Autonomous aerial capture of non-cooperative targets by Unmanned Aerial Vehicles (UAVs) introduces critical control and perception challenges caused by visual detection noise, sampling latency, temporary visual detection dropouts, and severe ground effect disturbances. To overcome these limitations within the context of retrieving entomological traps for the Vector Tracking System (VTS) project, this work develops and validates a closed-loop Position-Based Visual Servoing (PBVS) system utilizing a single camera. The proposed architecture structures the approach into three sequential operational phases, integrating specialized lightweight convolutional neural networks from the YOLO family (YOLO26n and YOLOv10n) tailored to each optical regime, achieving mean precision exceeding 99.2\% ($mAP@50$) with inference rates from 28.6 to 51.8 frames per second (fps). To attenuate the noise of the bounding boxes generated by YOLO, an Extended Kalman Filter (EKF) acts as a kinematic estimator, reducing fluctuations by 52\% to 91\% and ensuring trajectory continuity for up to 4.0 seconds of complete visual occlusion. Two-dimensional coordinates are projected into three-dimensional metric Cartesian space via an analytical Pinhole model, feeding a cascaded proportional-derivative (PD) flight controller with EKF-derived velocity to dampen oscillations and prevent overshoot. A dual-zone optical hysteresis strategy ensures 100\% field-of-view target retention, while attitude control converges with a mean yaw error of $0.46^\circ$ and a standard deviation of $\pm 0.54^\circ$. The system was designed to deploy a motorized robotic gripper mounted beneath the UAV airframe. Due to development timeline constraints, coupling feasibility was demonstrated in the physical simulator Gazebo Sim with ArduPilot in a Software-in-the-Loop (SITL) architecture, assessing the alignment of the gripper jaws relative to the trap handles under a $\pm 8\text{ mm}$ tolerance window. Across a stochastic campaign of 60 Monte Carlo missions with \textit{Ground Truth} validation, the system achieved a pre-drop alignment root mean square error (RMSE) of $3.13\text{ mm}$ at $0.20\text{ m}$ altitude, a terminal ground-contact RMSE of $5.25\text{ mm}$, and a 76.7\% coupling success rate, establishing a robust and millimetric solution for autonomous aerial manipulation without fiducial markers.
